\todo[inline]{How to structure the paper sections}
% \paragraph{Summary:}
% \paragraph{Problem:}
% \paragraph{Gap:}
% \paragraph{Method/Solution:}
% \paragraph{Findings:}
% \paragraph{Results:}
% \paragraph{Contribution:}
% \paragraph{Limitations:}
% \paragraph{Author Contribution:}
% \paragraph{Copyright Note: }
\paragraph{Summary (Introduction)}
\textbf{Purpose:} \\
To recap the context from previous chapters and highlight the key contributions of the current chapter.

\vspace{0.5em}
\textbf{It introduces:}
\begin{itemize}
    \item The gap in formal definitions
    \item The need to define "system adaptation" before "self-adaptive systems"
    \item The role of the MAPE-K model
    \item The proposed logical architecture
\end{itemize}

\vspace{0.5em}
\textbf{Use in your thesis:} \\
Begin each chapter with a brief overview of what was previously discussed and what this chapter contributes. This helps situate the reader.

%---------------------------------------------------------------
\paragraph{Problem Statement}
\textbf{Purpose:} \\
To clearly define the research problem(s) addressed in the chapter.

\vspace{0.5em}
\textbf{Problems include:}
\begin{itemize}
    \item Lack of clear definitions
    \item Conceptual confusion in terminology
    \item MAPE-K not differentiating between self-* systems
\end{itemize}

\vspace{0.5em}
\textbf{Use in your thesis:} \\
Clearly state the research gap and why it matters. This motivates your work and justifies your contribution.

%---------------------------------------------------------------
\paragraph{Gap Analysis}
\textbf{Purpose:} \\
To articulate what is missing in the current literature.

\vspace{0.5em}
\textbf{Identified gaps:}
\begin{itemize}
    \item Absence of a precise definition of "system adaptation"
    \item Lack of architectural blueprints
    \item Incomplete understanding of the adaptation process
\end{itemize}

\vspace{0.5em}
\textbf{Use in your thesis:} \\
Include a dedicated paragraph showing what others haven’t done yet, supported by citations. It strengthens the argument that your work is necessary.

%---------------------------------------------------------------
\paragraph{Method / Solution}
\textbf{Purpose:} \\
To describe what you are doing to solve the problem.

\vspace{0.5em}
\textbf{Proposed solution:}
\begin{itemize}
    \item Formal definition of system adaptation
    \item A logical architecture for CPSs
    \item A framework to bridge theory and implementation
\end{itemize}

\vspace{0.5em}
\textbf{Use in your thesis:} \\
Explain your proposed approach, framework, or method. Be specific and connect it to the identified problem.

%---------------------------------------------------------------
\paragraph{Findings / Key Questions for Framing}
\textbf{Purpose:} \\
To explain the foundational questions needed to define and analyze adaptivity.

\vspace{0.5em}
\textbf{Key framing questions:}
\begin{itemize}
    \item What is the function?
    \item What goals and context?
    \item What time scale and metrics?
\end{itemize}

\vspace{0.5em}
\textbf{Use in your thesis:} \\
Present any core framing questions or theoretical underpinnings that your method builds upon or introduces.

%---------------------------------------------------------------
\paragraph{Results}
\textbf{Purpose:} \\
To summarize the concrete results and contributions.

\vspace{0.5em}
\textbf{Your results might include:}
\begin{itemize}
    \item Validation of the formal definition
    \item Simulation setup and evaluation tools
    \item Generalizability of the Quality Function
\end{itemize}

\vspace{0.5em}
\textbf{Use in your thesis:} \\
Summarize the main outcomes and how they support your claims. Include both theoretical and empirical results.

%---------------------------------------------------------------
\paragraph{Contributions}
\textbf{Purpose:} \\
To make explicit what you contributed in this chapter.

\vspace{0.5em}
\textbf{Contributions include:}
\begin{itemize}
    \item Formalism
    \item Logical architecture
    \item Evaluation framework
    \item Metric for adaptation (Quality Function)
\end{itemize}

\vspace{0.5em}
\textbf{Use in your thesis:} \\
Have a clearly labeled paragraph listing your contributions. This helps reviewers and readers quickly see your impact.

%---------------------------------------------------------------
\paragraph{Limitations}
\textbf{Purpose:} \\
To acknowledge what your work doesn’t solve.

\vspace{0.5em}
\textbf{Limitations include:}
\begin{itemize}
    \item Ambiguities in the definition
    \item Partial observability in real CPSs
    \item Limitations of implementation
\end{itemize}

\vspace{0.5em}
\textbf{Use in your thesis:} \\
Be honest about limitations. This shows critical thinking and opens the door for future research.

%---------------------------------------------------------------
\paragraph{Author Contributions (optional)}
\textbf{Purpose:} \\
To credit each co-author for their role in collaborative work (common in paper-based theses).

\vspace{0.5em}
\textbf{Use in your thesis:} \\
Include if you are using published papers or co-authored chapters. This adds transparency.

\begin{itemize}
    \item \textbf{Author A}: Conceptual design, theoretical model.
    \item \textbf{Author B}: Implementation, validation.
    \item \textbf{Author C}: Experiments and analysis.
\end{itemize}

%---------------------------------------------------------------
\paragraph{Copyright Note (if applicable)}
\textbf{Purpose:} \\
To justify the inclusion of a reprinted article and show compliance with publisher rights.

\vspace{0.5em}
\textbf{Use in your thesis:} \\
Include if you are reprinting published articles. Follow your university’s thesis format.

\vspace{0.5em}
This chapter is based on the following published article: \\
\textit{Author names}, \textit{Paper title}, \textit{Journal Name}, \textbf{Volume}, Pages (Year). \\
DOI: \url{https://doi.org/xxxxxxx}. Reprinted with permission from the publisher.

% \todo[inline]{Example}

% \paragraph{Summary / Introduction}
% % Briefly summarize the previous chapter and introduce the content of this chapter.
% In this chapter, we build upon the findings from Chapter X and address two key challenges related to [...]. First, we tackle the lack of a formal definition of [...], and second, we propose an architectural framework for [...].

% \paragraph{Problem Statement}
% % State the main problems you aim to solve.
% Despite significant progress in [...], the engineering of [...] remains challenging. The absence of precise terminology and formal models hampers the development and evaluation of [...].

% \paragraph{Gap Analysis / Motivation}
% % Analyze the missing parts in the current literature.
% Existing literature proposes several models and frameworks for [...]; however, a precise, comprehensive, and widely accepted definition of [...] is still lacking. Additionally, architectural blueprints that support [...] in dynamic and uncertain contexts are rare.

% \paragraph{Proposed Method / Theoretical Contribution}
% % Present your solution or method.
% To address these issues, we propose a formal definition of [...], grounded in [...] theory. We also introduce a logical architecture for engineering [...] that supports [...].

% \paragraph{Theoretical Framework / Key Definitions}
% % Introduce the core definitions and theoretical framing.
% The foundation of our approach lies in defining the concept of [...]. We frame system adaptivity through the following key questions:
% \begin{itemize}
%     \item What is the system function (sf) considered adaptive?
%     \item What are the adaptation goals?
%     \item What context conditions trigger adaptation?
%     \item What time period is considered?
%     \item Under which convergence parameters is adaptivity assessed?
% \end{itemize}

% \paragraph{Results}
% % Describe what you achieved.
% We validate our approach through a model problem from the domain of [...]. We implemented a simulated system using [...], and evaluated the system adaptation using a metric referred to as the \textit{Quality Function}.

% \paragraph{Contributions}
% % Clearly state the novel contributions of this chapter.
% The main contributions of this chapter are:
% \begin{enumerate}
%     \item A formal definition of system adaptation.
%     \item A logical architecture for engineering [...]-adaptive systems.
%     \item A reusable evaluation framework with a Quality Function metric.
%     \item A differentiation between passive and active self-adaptation.
% \end{enumerate}

% \paragraph{Limitations}
% % Acknowledge what your work does not cover.
% While our theoretical contributions are validated, several limitations exist:
% \begin{itemize}
%     \item The definition of adaptation depends on partial observability.
%     \item Implementation only realizes passive self-adaptation.
%     \item Quality Function relies on full state observability, which may not generalize.
% \end{itemize}

% \paragraph*{Author Contributions (if applicable)}
% \textbf{Author 1}: Conceptualization, theoretical formalism, implementation.\\
% \textbf{Author 2}: Architecture design and review.\\
% \textbf{Author 3}: Evaluation framework and experiment execution.

% \paragraph*{Copyright Note (if applicable)}
% This chapter is based on a previously published article: \\
% \textit{Author Names}, \textit{Title}, \textit{Journal}, \textbf{Volume}, Pages (Year). \\
% DOI: \url{https://doi.org/xxxxxx}. Reprinted with permission from Elsevier.
